\section{Symmetry class of the sixteen-phase ensemble}
\label{sec:sixteen-phase-symmetry}
The canonical proof is
\path{mathematics/fixed-map/sixteen_phase_symmetry.md}.
The integer energy is certified by
\path{formalization/Bgs/SixteenPhaseSymmetryRational.lean}.
Python and C++ check the similarity and the momentum sine.
No external refereeing is asserted.
BGS-TRACE-005 and BGS-BGS-001 remain open.


\begin{theorem}[Symmetry class, BGS-TRACE-130]
\label{thm:sixteen-phase-symmetry}
Let \(A^{\mathrm{diag}}(\varphi)\) be the position-diagonal part of the sixteen-phase
symbol. For every real \(\eta\) and every phase vector, \(e^{i\eta A^{\mathrm{diag}}}M_N\)
is similar to the symmetric unitary \(D^{1/2}M_ND^{1/2}\). Every such realization
lies in the orthogonal class. A position sine does not leave that class.

The displacement-\(\ell=1\) modes have energy \(9/160000\). The momentum sine
\(B_N=(T_N(1,0)-T_N(-1,0))/(2i)\) is Hermitian and odd under position
conjugation, with symbol energy \(1/2\). For \(\varepsilon\ne 0\) the kick
\(A^{\mathrm{diag}}+\varepsilon B_N\) is not even under that conjugation, and
Theorem 1 does not apply. Neither statement identifies a COE or CUE form factor.
The time-one averaged squared trace remains the vanishing-coefficient statement
of BGS-TRACE-129.
\end{theorem}
